\begin{abstract}

零售销售明细同时存在强右尾、一单多行和仅48个月的短序列；若四问各用一套口径，订单计数、驱动判断、总量预测与资源分配将相互矛盾。本文沿“校正统计实体—筛选稳定关联—锁定预测模型—约束层级分配”推进，并以朴素基线、时间外验证和参数扰动检验各环节。

\textbf{对于问题1：}% src:交接/计划.json:问题清单[编号=1]
直接把明细行当订单会得到$9800$单，而按订单号去重仅有$4921$单，高估$99.15\%$。为此构建“订单—月份双口径稳健画像”：原尺度保留销售总额，订单层用归并与分位数刻画典型交易，月份层识别季节结构。结果显示，最高一成明细贡献$60.40\%$的销售额，第四季度占$38.52\%$，销售兼具头部集中与年末后置特征。% src:求解/问题1/结果/验证与灵敏度.json:朴素基线;求解/问题1/结果/核心统计与业务指标.json:明细层Sales分布.Top10%样本销售额占比;求解/问题1/结果/时间结构.json:季度季节结构[季=Q4].销售额份额

\textbf{对于问题2：}% src:交接/计划.json:问题清单[编号=2]
一张订单可能包含多个品类，整单汇总会丢失产品信息，逐行检验又会制造伪独立。因此提出“订单簇置换的双证据因子筛”，在明细层保留产品结构，但按订单整簇完成检验、重采样与损失计算，并用跨年方向和参数网格复核。产品历史价值层、产品大类和产品子类在$81/81$组扰动中均入选；区域仅入选$54/81$组，故降为次级信号。同一2025年时间外样本上，订单等权log1p MAE由$1.105920$降至$0.590432$，改善$46.61\%$。% src:求解/问题2/结果/关键参数灵敏度.json:各因素入选次数,组合总数;求解/问题2/结果/基线对比.json:计划内同口径朴素基线

\textbf{对于问题3：}% src:交接/计划.json:问题清单[编号=3]
训练期仅36个月、测试期仅12个月，不宜在封存结果上反复择模。本文在训练期锁定“季节锚定—偏差校正组合预测”，以季节朴素保留去年同月结构，用阻尼ETS与谐波岭修正水平、趋势和季节形状，并校正对数反变换偏差。2025年封存测试中，组合相对季节朴素将MSE、MAE和MAPE分别降低$60.42\%$、$37.95\%$和$34.17\%$，四个季度MAE均胜出，九组权重扰动下部署结论不变。% src:求解/问题3/结果/数据切分.json;求解/问题3/结果/2025评估指标.json;求解/问题3/结果/组合权重与灵敏度.json:关键参数灵敏度

\textbf{对于问题4：}% src:交接/计划.json:问题清单[编号=4]
各层级独立预测会造成产品大类与子类无法回加，故采用“总量锚定的层级份额收缩”：按稳定性融合长期、近期与同月份额，再以父级份额乘条件子类份额保持协调。2025年五维平均份额WAPE由固定份额基线的$39.585\%$降至$29.015\%$，各维度最大加总误差不超过$5.82\times10^{-11}$。2026年销售额点预测为$791951.82$，全年可加总风险边界为$[454803.67,1379541.60]$；2025年封存测试的$95\%$经验区间覆盖率仅为$83.33\%$，故该边界不作严格概率保证，Technology、Phones和Storage可优先配置，但仍需保留悲观情景下的规模防守。% src:求解/问题4/结果/份额收缩回测与灵敏度.json:2025五维平均对比;求解/问题4/结果/2026总量与一致性.json:层级一致性检查,2026全年总量预测;求解/问题3/结果/2025评估指标.json:区间测试表现.95%覆盖率;交接/结果解读_问题3.md:局限与使用边界;求解/问题4/结果/机会_把握度策略矩阵.json:点预测与下界排序稳健性

四问依次回答“数什么、信什么、怎样预测、如何分配”，把数据口径、模型精度和策略强度约束在同一证据链内。

\keywords{稳健统计；订单簇置换；组合预测；层级份额收缩；零售经营}
\label{abstract:end}

\end{abstract}
